\begin{textsample}{NEG Dim 2 – global\_south\_2024\_12 – Score 1.00 – global\_south\_2024\_12/118655718.txt}  \label{ex:f2_neg_383}
South African research funds ' at risk ' over Gaza stance A number of local universities may face cuts to their funding after cutting ties with Israeli institutions. ( Photo by Julian Stratenschulte/picture alliance via Getty Images ) Solidarity with Palestine in the ongoing conflict in the Middle East is threatening research funding at some South African universities, according to former cabinet minister Naledi Pandor.
The former international relations and cooperation minister, who has also served as minister of science and technology, said several universities could lose funding after severing ties with Israeli institutions and urged the South African government to offer financial support.
Pandor told the Science Forum South Africa in Pretoria last week that the sector was not immune to the impact of the toxic geopolitical environment."South Africa ' s support for the struggle for Palestinian sovereignty has resulted in some key research initiatives that rely on international funding facing the threat of funds being withdrawn,"she said."This is reportedly happening to universities that have decided not to pursue links with institutions in Israel that have links to the military actions in Palestine."Pandor, who retired in June, said the department of science, technology and innovation must give attention to the impact of the withdrawal of funds on institutions and researchers and support them in identifying alternative resources.
The department co-hosted the Science Forum with the Science Diplomacy Capital for Africa initiative, under the theme Igniting Conversations about Science -- Innovation and Science for Humanity.
Pandor called for the event to be more than"a mere talk shop", urging delegates :"It must play a full role in advancing African capabilities and ensuring that Africa rising becomes a reality and not a populist slogan."Universities South Africa, the country ' s representative body for 26 vice-chancellors at public institutions, confirmed that at least one university had experienced funding cuts over its stance on Gaza, without elaborating.
It declined to comment further, saying the matter still had to be discussed among vice-chancellors and that each institution could speak for itself on its position.
Salim Vally, a professor in the Faculty of Education at the University of Johannesburg, said that any donors threatening to withdraw funds were indulging in academic blackmail."Universities must not be deterred and should remain true to their mission, which must be an ethical commitment to seeking the truth, social justice, human rights, anti-racism, solidarity and knowledge that benefits humanity,"he said."This is one of the ideals of knowledge for the public good."Vally believes Israel ' s assault on Gaza, where all universities have been bombed and more than 12 000 students, plus scores of faculty members, have been killed, is a litmus test for all intellectuals, academics and university management.
He described as"cowardice"the refusal of some universities to take a stand against Israel and applauded Pandor ' s attempts to mitigate any loss of funding from those that do."Taking a principled stand often comes with sacrifice,"added Vally, who is also a prominent human rights activist."South Africans who fought against our erstwhile apartheid regime and who called on the world to boycott apartheid state institutions know this well."Israel has faced an unprecedented academic boycott since it launched its war on Gaza in October last year, which has killed more than 44 500 Palestinians.
A growing number of European universities are among those taking action.
Earlier this year, the International Science Council ( ISC ) updated its position on academic boycotts, described as a collective protest by an academic community or institution to express disapproval of other academics or institutions, or to put pressure on them to meet demands."The ISC, as a general principle, does not endorse academic boycotts,"the statement said, citing article 27 of the Universal Declaration of Human Rights and its own Principles of Freedom and Responsibility in Science."Exceptions to the general principle will be considered by the ISC governing board when there are clear and systemic violations of human rights,"it added.
To provide the best experiences, we use technologies like cookies to store and/or access device information.
Consenting to these technologies will allow us to process data such as browsing behavior or unique IDs on this site.
Not consenting or withdrawing consent, may adversely affect certain features and functions.
FunctionalFunctional Always active The technical storage or access is strictly necessary for the legitimate purpose of enabling the use of a specific service explicitly requested by the subscriber or user, or for the sole purpose of carrying out the transmission of a communication over an electronic communications network.
PreferencesPreferences The technical storage or access is necessary for the legitimate purpose of storing preferences that are not requested by the subscriber or user.
StatisticsStatistics The technical storage or access that is used exclusively for statistical purposes.The technical storage or access that is used exclusively for anonymous statistical purposes.
Without a subpoena, voluntary compliance on the part of your Internet Service Provider, or additional records from a third party, information stored or retrieved for this purpose alone can not usually be used to identify you.
MarketingMarketing The technical storage or access is required to create user profiles to send advertising, or to track the user on a website or across several websites for similar marketing purposes.

% matched lemmas: 
\end{textsample}
